\section{Service design}
\label{sec:design}

\subsection{Architecture}

The service is a single pygeoapi instance exposing two OGC API families over
two data stores. Figure~\ref{fig:arch} shows the whole system, including where
the data comes from and which part of it stands in for Bhoonidhi.

\begin{figure}[htbp]
\centering
% Scaled to the measure: the diagram's natural width depends on node
% placement, so resizebox guarantees it never spills into the margin.
\resizebox{\textwidth}{!}{%
\begin{tikzpicture}[
  font=\sffamily\scriptsize,
  node distance=5mm,
  box/.style={draw=rule,fill=white,rounded corners=2pt,align=center,
              inner sep=4pt,minimum height=8mm},
  upstream/.style={box,fill=shade,draw=rule,text width=72mm},
  stage/.style={box,text width=72mm},
  tier/.style={box,draw=accent,line width=0.7pt,text width=52mm},
  proc/.style={box,draw=rule,text width=46mm,font=\sffamily\tiny},
  lbl/.style={font=\sffamily\tiny\itshape,color=comment,align=left},
  ar/.style={-{Stealth[length=2mm]},draw=comment}
]

% upstream
\node[upstream] (bh) {%
  \textbf{Bhoonidhi Open Data Access API} (ISRO/NRSC)\\
  STAC 1.0.0: \texttt{/auth/token} $\rightarrow$ \texttt{/data/search} $\rightarrow$ \texttt{/download}\\
  \textcolor{bad}{spatial surface = \texttt{bbox} + \texttt{intersects}; computes nothing}};

% harvest
\node[stage,below=7mm of bh] (harv) {%
  \texttt{scripts/fetch\_footprints.py} \quad\textbar\quad \texttt{build\_catalogue.py}\\
  STAC client: \texttt{--source bhoonidhi} (JWT) or \texttt{--source open}};

\draw[ar] (bh) -- (harv);

% stores
\node[box,below left=7mm and -22mm of harv,text width=42mm] (fps) {%
  \textbf{footprints.geojson}\\
  68 real acquisition footprints\\
  \textcolor{comment}{60 optical + 8 radar}};
\node[box,below right=7mm and -22mm of harv,text width=42mm] (cat) {%
  \textbf{records.tinydb}\\
  59 collection records\\
  \textcolor{comment}{provenance-tagged}};

\draw[ar] (harv) -- (fps);
\draw[ar] (harv) -- (cat);

% dem
\node[box,left=6mm of fps,text width=22mm] (dem) {%
  \textbf{dem.tif}\\
  \textcolor{comment}{raster}};

% service
\node[tier,below=9mm of fps] (proc) {%
  \textbf{OGC API --- Processes}\\
  \textcolor{comment}{18-062r2 Core}\\
  sync + async, job API};
\node[tier,below=9mm of cat] (rec) {%
  \textbf{OGC API --- Records}\\
  \textcolor{comment}{20-004, supporting}\\
  full-text, queryables, bbox};

\draw[ar] (fps) -- (proc);
\draw[ar] (dem.south) |- ([yshift=4mm]proc.north west) -- (proc.north west);
\draw[ar] (cat) -- (rec);

\begin{scope}[on background layer]
\node[draw=accent,dashed,rounded corners=3pt,inner sep=5mm,
      fit=(proc)(rec),label={[font=\sffamily\tiny\bfseries,color=accent]below:%
      pygeoapi 0.24.0 \textperiodcentered\ one instance \textperiodcentered\ port 5003}] (svc) {};
\end{scope}

% processes
\node[proc,below=13mm of svc.south,text width=92mm] (plist) {%
  \textbf{footprint geometry} \quad \texttt{aoi-coverage} \textperiodcentered\
  \texttt{minimal-scene-set} \textperiodcentered\ \texttt{multi-sensor-overlap}
  \textperiodcentered\ \texttt{revisit-density}\\[2pt]
  \textbf{raster values} \quad \texttt{zonal-statistics} \textperiodcentered\
  \texttt{scene-raster-profile} \textcolor{comment}{(the join)}};

\draw[ar] (svc.south) -- (plist.north);

\end{tikzpicture}}
\caption{Service architecture. Footprints are harvested from a STAC API ---
Bhoonidhi's own with credentials, or an open substitute without --- and the
Processes tier computes over that geometry and the raster beneath it. The
Records tier is supporting infrastructure, giving the processes discoverable
collection identifiers to point at.}
\label{fig:arch}
\end{figure}

\subsection{The two layers, deliberately}

The six processes divide cleanly, and the division is the design:

\begin{itemize}
  \item the \textbf{footprint processes} read only geometry. They answer
        \emph{which scenes are worth fetching} --- and they answer it
        \emph{before} anything is downloaded, which is the entire point;
  \item the \textbf{raster processes} read pixels. They answer \emph{what is
        actually inside} a polygon;
  \item \code{scene-raster-profile} joins the two: for every scene footprint
        intersecting the AOI it clips the raster to that footprint and reports
        what lies underneath, ranked by overlap. Choosing between candidate
        scenes on what they actually contain is impossible from a catalogue
        that only answers \emph{does this intersect your box}.
\end{itemize}

This is deliberately \emph{not} spectral index computation. There is no band
maths and no derived imagery product, because the gap identified in
Section~\ref{sec:gap} is spatial algebra over the catalogue's own geometry,
and a service should fill the gap it claims to fill.

\subsection{Why these six operations}

Each maps directly onto a row of Table~\ref{tab:questions} --- an ordinary
question that Bhoonidhi's single predicate cannot answer. The mapping is
one-to-one by construction; the processes were chosen from the audit, not the
other way round.

\code{minimal-scene-set} is the sharpest of them and deserves its own note.
It solves a \textbf{set cover} problem: given every footprint touching the
AOI, find the fewest whose union covers it. The implementation is greedy ---
at each step it takes the footprint adding the largest previously-uncovered
area, stopping when the target fraction is reached or no candidate adds more
than 1000\,m\textsuperscript{2} (which is noise, not coverage).

Greedy set cover is a deliberate choice over an exact solver. Set cover is
NP-hard~\cite{setcover}; the greedy algorithm is within a $\ln(n)$ factor of
optimal, runs in time acceptable for a synchronous HTTP request, and produces
a \emph{cumulative coverage curve} that is more useful to a user than a bare
optimal set --- it shows what each additional scene buys. The response states
the algorithm and its approximation bound in a field named \code{algorithm},
so no client can mistake the result for a proven minimum.

\subsection{Measurement: areas in an equal-area projection}

Every area in every response is computed in a \textbf{Lambert Azimuthal
Equal-Area} projection centred on \emph{that request's own} area of interest:

\begin{lstlisting}
+proj=laea +lat_0=<aoi centroid lat> +lon_0=<aoi centroid lon>
+datum=WGS84 +units=m +no_defs
\end{lstlisting}

This matters more than it may appear. The whole service exists to answer
area-based questions --- \emph{how much is covered}, \emph{what is the gap},
\emph{which scene adds the most}. Computing those in degrees would make every
number meaningless, and silently so: a square degree is neither a constant
area nor an area at all. Centring the projection per request rather than using
one national CRS keeps distortion negligible across an AOI of any size
anywhere in India, at the cost of one \code{pyproj} transformer construction
per call.

A compliance test asserts that a $0.5^\circ \times 0.7^\circ$ AOI reports an
area in the thousands of km\textsuperscript{2} rather than a fraction, so this
cannot silently regress.

\subsection{Honesty about provenance, in the response}

Every process response carries a \code{provenance} block stating how many
footprints were considered, which collection and window, and --- critically
--- whether the footprint store is Bhoonidhi's or a substitute. The store
metadata carries \code{is\_substitute\_for\_bhoonidhi: true} and this
propagates into every single response.

This is a design decision, not a disclaimer bolted on. A spatial analysis
service whose inputs are unclear is worse than none, because its outputs look
authoritative. No result this service returns can be mistaken for Bhoonidhi
output.

\subsection{Implementation choices, and why}

\begin{center}
\small
\begin{tabularx}{\textwidth}{@{}>{\raggedright\arraybackslash}p{0.27\textwidth}>{\raggedright\arraybackslash}X@{}}
\toprule
\textbf{Decision} & \textbf{Reasoning} \\
\midrule
\textbf{OGC API --- Processes}, not WPS 1.0.0/2.0
  & The modern successor: JSON over REST, a crisp mandatory operation set, a
    real job model, and an OpenAPI description. It demonstrates far better
    than XML, and it is the standard a new service in 2026 should implement. \\
\textbf{pygeoapi}
  & Reference implementation of the OGC API family, so conformance comes from
    the framework rather than from hand-rolled routing. Compliance is the
    largest block of marks and hand-rolling a job API to spec is the wrong
    risk to take. \\
\textbf{shapely + pyproj}
  & The geometry and projection engines the rest of the Python geospatial
    stack is built on. The processes are thin: the hard parts are delegated
    to libraries that are correct. \\
\textbf{Greedy set cover}
  & See above: NP-hard exactly, within $\ln(n)$ greedily, and the coverage
    curve is more informative than a bare optimum. Declared in the response. \\
\textbf{Per-request LAEA}
  & Area-based questions need real areas; one national CRS would distort at
    the extremes. \\
\textbf{TinyDB job manager}
  & Jobs are persisted, so asynchronous execution survives a restart and
    \ep{/jobs/\{id\}} keeps working --- which OGC API --- Processes expects.
    A production deployment would use PostgreSQL. \\
\textbf{Jobs under \code{runtime/}}
  & Originally \code{/tmp}, which the OS clears on reboot; a missing output
    directory surfaces as an opaque HTTP 400 on \emph{every} execution.
    See Section~\ref{sec:defects}. \\
\textbf{Native deployment as the tested path}
  & The build host's user was not in the \code{docker} group and no sudo
    password was available. The native path is the verified one and is
    labelled as such rather than presenting an unexercised container path as
    primary. \\
\bottomrule
\end{tabularx}
\end{center}
